\chapter{Discusión de resultados}
Los resultados permiten analizar el funcionamiento del sistema desde
tres niveles relacionados entre sí. Estos son, la capacidad de registrar
el comportamiento producido durante las interacciones, la representación
de ese comportamiento mediante estados y la modificación de las condiciones
de juego por medio del sistema de dificultad dinámica adaptaivo 
(DDA, por sus siglas en inglés). Las respuestas de la encuesta complementan esta 
información al proporcionar una valoración subjetiva posterior a la 
experiencia. Esta combinación permite evaluar el ciclo de medición,
interpretación y adaptación planteado en el objetivo general sin asumir 
que las métricas objetivas representan directamente la percepción del 
participante.

\section{Desempeño registrado y capacidad de medición}
Las diferencias observadas en el cuadro \ref{tab:metricas} muestran que el
sistema registra comportamientos distintos entre participantes y entre mecánicas.
Al considerar conjuntamente los intentos de cada minijuego, las tasas de
éxito se distribuyen desde, aproximadamente, un 35.37\% en Juan Andrés hasta
un 71.09\% en Joaquin, mientras que los demás se encuentran entre estos extremos.
Esta variación resulta relevante para el sistema adaptativo, ya que confirma
que una configuración uniforme de dificultad no produce el mismo nivel de
desempeño en todas las actividades ni para todos los participantes.

Las métricas específicas complementan esta información al caracterizar
aspectos de la ejecución que la tasa de éxito por sí sola no permite
distinguir. El cuadro \ref{tab:precisionManoPromedio} muestra diferencias
entre la dirección realizada y el uso de la mano requerida en Shipi, Juan
Andrés y Jeremy, mientras que el cuadro \ref{tab:reaccionPromedio} presenta
tiempos de reacción que varían entre 0.25 y 0.92 segundos entre participantes.
Esta separación resulta consistente con el uso de múltiples indicadores 
de comportamiento en el modelado del jugador, ya que un único resultado
de éxito o fallo no describe completamente la forma en que se ejecuta
una acción.

El cuadro \ref{tab:timingPromedio} también evidencia por qué las métricas
específicas deben interpretarse dentro del contexto de cada mecánica. Los
errores temporales presentan escalas diferentes entre minijuegos porque el
instante esperado y la ventana disponible para actuar no representan
exactamente la misma acción en todas las actividades. Por esta razón, un
error de \textit{timing} mayor en una mecánica no implica necesariamente
un peor desempeño que un valor menor en otra, y el análisis utiliza estas
variables principalmente para caracterizar el comportamiento dentro de cada
actividad.

En conjunto, estos resultados respaldan el mecanismo de captura definido
en el segundo objetivo específico \ref{sec:objetivosEspecificos}. Los cuatro
cuadros ateriormente mencionados contienen información procedente de los seis
minijuegos y de todos los participantes, lo que permite reconstruir tanto
el resultado general de una acción como características específicas de
su ejecución. La telemetría no queda limitada a una puntuación final, sino que
conserva información suficiente para analizar posteriormente tendencias y
diferencias de desempeño, condición necesaria para alimentar un sistema de
adaptación basado en comportamiento.

\section{Modelado de estados y comportamiento del DDA}
En el cuadro \ref{tab:cantidadEstados} se muestra que el modelo no permanece
en una única clasificación durante las sesiones. En los registros aparecen 
225 entradas a \textit{idle}, 357 a \textit{engaged} y 212 a \textit{overloaded},
lo que indica que las reglas implementadas permiten representar variaciones del
comportamiento durante la experiencia. Estos valores corresponden a transiciones
registradas y no al tiempo total que el participante permanece en cada estado,
por lo que su interpretación se limita a la frecuencia con la que el modelo cambia
su clasificación.

La presencia de los tres estados operativos entre los participantes resulta
coherente con el propósito del modelado de sintetizar tendencias recientes
en lugar de representar una característica permanente del usuario. Los modelos
temporales de jugador consideran que el desempeño puede cambiar conforme el usuarioa
aprende una mecánica, encuentra nuevas condiciones o presenta variaciones durante 
una misma sesión. En este caso el uso de una ventana reciente permite que una 
misma persona pase entre períodos de actividad regular, inactividad y acumulación
de errores sin mantener una clasificación fija para toda la experiencia.

Esta variabilidad se refleja en las decisiones del DDA. El cuadro 
\ref{tab:cambiosPorDificultad} contiene un total de 216 modificaciones entre los 34
participantes, de las cuales 96 corresponden a aumentos y 120 a 
disminuciones, y todos los participantes presentan cambios en algún
momento de la experiencia. En conjunto con el cuadro 
\ref{tab:cambiosDificultad}, esto muestra que el sistema no sigue una progresión
predeterminada hacia dificultades mayores, sino que se desplaza en ambas direcciones
de acuerdo con el desempeño reciente y el estado identificado.

El cuadro \ref{tab:maxDificultad} refuerza esta característica al mostrar
que los participantes no alcanzan los mismos niveles máximos en todos los minijuegos.
Parte de esta diferencia responde a las características particulares de 
cada mecánica, pero también debe considerarse que el nivel global se 
conserva entre minijuegos, por lo que el comportamiento registrado en segmentos
anteriores influye en el nivel con el que se inicia una actividad posterior.

La evidencia más directa sobre la respuesta del sistema se encuentra en 
la figura \ref{fig:rendimientoDDA}, donde se compara el desempeño
inmediato antes y después de los ajustes. En las 68 reducciones de dificultad
que cumplen las condiciones para realizar esta comparación, la tasa de
éxito pasa en promedio de 34.71\% a 62.06\%, lo que corresponde a un aumento de 27.35
puntos porcentiales. 53 de los 68 casos presentan una mejora, 12 premanecen
sin cambio y 3 disminuyen. En sentido contrario, los 69 aumentos analizables
parten de una tasa promedio de 91.30\% y luego presentan 60.00\% con una disminución 
media de 31.30 puntos porcentuales. 52 casos reducen de su desempeño, 16 permanecen sin cambio
permanecen sin cambio y únicamente uno muestra un aumento.

La dirección de estos cambios es consistente con la función asignada al 
DDA. Una reducción ocurre después de identificar dificultades recientes 
y es seguida frecuentemente por una recuperación de la tasa de éxito, 
mientras que un incremento ocurre después de un período de rendimiento 
alto y produce una condición de mayor exigencia en la que el porcentaje 
de aciertos disminuye. Esto aporta evidencia de que los niveles fácil, 
normal y difícil no funcionan únicamente como etiquetas internas, 
sino que las modificaciones de velocidad, cantidad y complejidad 
asociadas con ellos producen diferencias observables en la ejecución, 
tal como se espera de un sistema de ajuste dinámico de dificultad.

Sin embargo, esta comparación no permite atribuir causalmente toda la 
variación al cambio de dificultad. Los incrementos se activan 
precisamente después de ventanas con un desempeño alto y las reducciones 
después de condiciones de bajo desempeño o sobrecarga, por lo que una 
parte de la diferencia posterior puede estar asociada con la propia 
forma en que se seleccionan los momentos de adaptación. Aun con esta 
limitación, la consistencia en la dirección de los cambios muestra que el 
sistema responde a variaciones del jugador y modifica las condiciones de 
interacción en la dirección prevista, mientras que las ventanas de 
observación evitan utilizar un único evento aislado como criterio de 
adaptación.

\section{Correspondencia entre telemetría y percepción}
La comparación entre las métricas objetivas y las respuestas posteriores 
muestra que ambas fuentes describen dimensiones diferentes de la 
experiencia. En la figura \ref{fig:frustracionVSExito}, el minijuego 
seleccionado como más frustrante coincide con el de menor tasa de éxito 
únicamente para 13 de los 34 participantes. Aunque Juan Andrés aparece 
tanto como el minijuego con la menor tasa de éxito global como el 
señalado con mayor frecuencia como más frustrante en la figura
\ref{fig:juegoFrustrante}, esta relación general no se reproduce de forma
consistente a nivel individual.

Este resultado indica que la frustración percibida no puede 
representarse únicamente mediante la proporción de aciertos. Factores 
como el tipo de movimiento requerido, la cantidad de condiciones que 
deben cumplirse simultáneamente, la claridad de las indicaciones o la 
familiaridad con una mecánica pueden intervenir en la valoración del 
jugador, aunque estas variables no se miden de forma independiente en 
este estudio. La diferencia entre ambos tipos de información coincide 
con la necesidad de combinar datos objetivos y subjetivos al evaluar 
experiencias de realidad virtual (VR, por sus siglas en inglés), debido 
a que la telemetría describe el 
comportamiento observable pero no sustituye directamente la valoración 
realizada por el participante.

Una situación similar aparece al comparar la valoración del desempeño 
con la tasa de éxito registrada. Como se observa en la figura
\ref{fig:desempenioVSExito}, los participantes que califican su 
desempeño como malo presentan una tasa media cercana al 56.27\%, 
quienes lo consideran intermedio alcanzan aproximadamente 59.38\% y 
quienes lo califican como bueno presentan aproximadamente 56.48\%. La 
ausencia de una progresión clara entre estas categorías muestra que 
la percepción global del propio desempeño no corresponde de forma 
directa con la proporción total de interacciones exitosas.

La comparación del agobio percibido con el estado \textit{overloaded} 
muestra una correspondencia igualmente limitada. La figura \ref{fig:FrecuenciaOverloaded}
presenta una amplia superposición entre las categorías de agobio y la 
frecuencia de entradas a \textit{overloaded} por cada 100 interacciones, 
sin observarse un incremento progresivo a medida que aumenta la 
valoración subjetiva. Los promedios se mantienen cercanos entre las 
categorías, por lo que el estado computacional no funciona como una 
medida directa de cuánto declara sentirse abrumado el participante.

Este resultado no contradice la función definida para \textit{overloaded}, 
ya que el estado se construye a partir de errores y omisiones recientes y 
no como una estimación psicológica del usuario. Su utilidad se encuentra 
en identificar patrones de desempeño que justifican una respuesta del DDA, 
algo que se observa posteriormente en la recuperación asociada con las 
reducciones de dificultad de la figura \ref{fig:rendimientoDDA}. Por lo
tanto, los resultados apoyan el uso de \textit{overloaded} como un 
estado operacional para la adaptación, pero no respaldan utilizarlo como 
equivalente del agobio subjetivamente reportado.

\section{Relación entre exigencia de interacción y atención a la música}
La atención a la canción constituye un indicador relevante debido a que 
las mecánicas se incorporan dentro de un videoclip gamificado y no como 
actividades aisladas. La figura \ref{fig:atencionCancion} muestra que 23
de 34 participantes, equivalente al 67.65\%, seleccionan valores de 4 o 5
en la escala de atención, y el promedio general alcanza 3.74 sobre 5. Al
mismo tiempo, la figura \ref{fig:diversion} muestra que 32 de las 33 
respuestas válidas califican la diversión con 4 o 5, con un promedio de 
4.42, lo que indica que la presencia de mecánicas interactivas resulta 
compatible con una valoración alta de entretenimiento durante la 
experiencia.

La percepción de carga tampoco se concentra en los niveles superiores. 
Como muestra la figura \ref{fig:nivelAgobio}, 25 de los 34 participantes 
se describen como nada o ligeramente abrumados, mientras que solamente 
uno selecciona la categoría de mayor agobio. Considerados junto con la 
atención reportada, estos resultados muestran que la interacción puede 
coexistir con la escucha de la canción para una proporción mayoritaria 
de la muestra, aspecto importante en una experiencia donde el 
\textit{gameplay} debe acompañar la música sin sustituirla completamente.

No obstante, la figura \ref{fig:atencionVSDDA} no presenta una 
relación clara entre la atención y la cantidad de estados 
\textit{overloaded} o reducciones de dificultad. Los participantes con 
niveles altos de atención aparecen tanto en sesiones con pocas como con 
varias intervenciones del sistema, y el mismo solapamiento se observa en 
los niveles inferiores. En consecuencia, los datos permiten afirmar que 
la experiencia adaptativa mantiene niveles favorables de atención musical 
en buena parte de la muestra, pero no permiten atribuir esa atención 
específicamente al funcionamiento del DDA.

\section{Limitaciones}
La evaluación presenta limitaciones que delimitan el alcance de estas 
conclusiones. De los 216 cambios de dificultad registrados en la tabla
\ref{tab:cambiosDificultad}, 137 disponen de cinco interacciones válidas 
antes y después bajo las condiciones necesarias para realizar la 
comparación de la figura \ref{fig:rendimientoDDA}, por lo que el análisis 
de variación inmediata no representa todos los ajustes realizados. 
Además, las métricas específicas no son equivalentes entre minijuegos, 
el estado \textit{overloaded} representa una clasificación operacional y 
no una medida subjetiva, y la ausencia de una condición sin DDA impide 
atribuir causalmente al sistema las valoraciones de atención, diversión 
o agobio.

Finalmente, el nivel de dificultad global se conserva a lo largo de una 
secuencia fija de minijuegos, lo que introduce una dependencia entre el 
desempeño previo y las condiciones encontradas posteriormente. Este 
diseño permite mantener continuidad en la adaptación, pero dificulta 
separar completamente el efecto de la mecánica del efecto acumulado de 
segmentos anteriores al comparar la tabla \ref{tab:maxDificultad}. Por
ello, los resultados sustentan el funcionamiento del sistema 
implementado y su capacidad de responder al comportamiento registrado, 
mientras que una evaluación futura con orden contrabalanceado o 
condiciones con dificultad fija permitiría estudiar con mayor precisión 
la efectividad de la adaptación sobre la experiencia del usuario.